\documentclass[instructions.tex]{subfiles}
\begin{document}
 
\chapter{CHAPITRE LXXIII. DU SACREMENT DE PÉNITENCE. Avis pour la confession.}
 
\primo Il y a de la différence entre la vertu de pénitence, et le sacrement de pénitence. La vertu de pénitence, dont nous avons parlé, nous est nécessaire pour conserver la grace sanctifiante : elle nous dispose à recevoir le pardon du péché, et l’expie quand il est pardonné; mais c’est dans le sacrement de pénitence que nous recevons le pardon, et la grace quand nous l’avons perdue. Nos pénitences sont les actions et les mérites de l’homme, qui sont peu de chose, si on les sépare de la grace; mais dans le sacrement de pénitence, ce sont les mérites de Jésus-Christ qui nous sont appliqués, et qui effacent, en un moment, ce que les mérites de tous les hommes ne pourraient effacer.
 
Tel est le privilége de la loi nouvelle, dans laquelle il est bien plus facile de recouvrer l’amitié de Dieu, et de recevoir le pardon, qu’il ne l’était dans l’ancienne loi; privilége divin accordé par Jésus-Christ à son Église. Le Sauveur, donnant la mission à ses apôtres pour établir cette Église sainte, leur dit : Recevez le Saint-Esprit : ceux dont vous remettrez les péchés, ils seront remis; et ceux dont vous les retiendrez, ils seront retenus\footnote{(1) Joan. 20. 22. 23.}. Pouvoir auguste, qui a toujours été exercé dans l’Église de Jésus-Christ, et auquel se sont soumis les plus grands génies, les plus grands princes de l’univers, les nations entières; pouvoir que les hérétiques ont combattu par malice, et auquel ils se sont soustraits par aveuglement; pouvoir consolant pour les fideles. Lorsqu’on s’avoue coupable au tribunal des hommes, on y est condamné; mais au tribunal de la pénitence, aussitôt qu’on s’avoue criminel avec un repentir sincère, on y est absous et justifié.
 
\secundo Que nous serions coupables de négliger ce sacrement, et d’abuser, pour notre perte, d’un remède que Dieu a établi pour notre sanctification! Allons donc souvent nous purifier dans ce bain sacré de Jésus-Christ. Nous sommes foibles, exposés aux tentations, et nous tombons souvent; allons-y chercher le remède, la guérison et la force; déclarons-y nos péchés entierement et avec sincérité. De quoi sert le déguisement devant un Dieu qui voit les plus secrets replis de notre cœur? Déclarons-les avec simplicité; faisons-en connoître le nombre et les circonstances nécessaires, sans embarrasser nos confessions de tant de paroles, de tant de récits, de tant de minuties et de circonstances inutiles. Accusons nos péchés avec confusion, nous regardant comme des criminels de lése-majesté divine, qui ne méritons que l’enfer. Accusons-les avec confiance, et avec amour pour un Dieu qui est prêt à nous recevoir dans le sein de sa miséricorde, aussitôt qu’il nous voit gémir à ses pieds. Accusons-nous enfin dans un esprit de religion et de foi; respectons dans le confesseur l’autorité de Jésus-Christ, ou plutôt ne regardons que Jésus-Christ dans son ministere, et n’oublions pas que l’accusation doit être accompagnée d’un propos efficace de se corriger.
 
\section{Résolutions}
 
\primo J’aurai recours au sacrement de pénitence souvent, et au moins toutes les fois que je me sentirai en danger de tomber dans le péché mortel, par le relâchement qui survient peu à peu, à mesure qu’on s’éloigne de ce salutaire remede.
 
\secundo J’y aurai recours aussi, dès que j’aurai fait quelque faute grave.
 
\tertio Mais j’y porterai toujours les sentimens de sincérité, d’humilité, de repentir et de ferme propos que ce sacrement exige.
 
\prayer{Mon Dieu, faites que je ne néglige désormais plus ce moyen efficace, que votre miséricorde nous offre pour guérir les plaies que le péché a faites à notre âme; et ne permettez pas que jamais j’en abuse.}
 
\end{document}
